% Mathematical Addendum — MATH-REGISTER VERSION (de-jargoned, for Prof. Fierro)
% Companion to addendum.tex. Same mathematics, same citations, same numbers —
% but computational/systems vocabulary is translated to its mathematical
% equivalent, results are stated as formal propositions, and empirical
% corroboration is relegated to one neutral sentence per section.
% Produced 2026-07-13 (Opus) per rmw3's calibration: "let the mathematics flower."
\documentclass[11pt]{article}
\usepackage[margin=1.1in]{geometry}
\usepackage{amsmath, amssymb, amsthm}
\usepackage[T1]{fontenc}

\newtheorem{proposition}{Proposition}
\newtheorem{remark}{Remark}

\newcommand{\R}{\mathbb{R}}
\newcommand{\E}{\mathbb{E}}
\newcommand{\Om}{\Omega}

\title{Mathematical Addendum: Existence, Convergence, and Stochastic\\
Viability of the Feasible Region $\Omega$}
\author{VSF Research Group}
\date{July 2026}

\begin{document}
\maketitle

\section*{Context}

We are grateful to the reviewer for three precise objections to the original
existence-and-invariance argument for the feasible region $\Omega \subset \R^{n}$
($n=128$) under the stochastic recurrence
\begin{equation}
  x_{t+1} = f_0(x_t) + \sigma\,\eta_t,
  \qquad t = 0, 1, 2, \dots
  \label{eq:rec}
\end{equation}
where $f_0 : \R^n \to \R^n$ is a continuous map, $(\eta_t)$ are i.i.d.\ zero-mean
perturbations, and $\sigma > 0$. The objections concerned (i) the convexity of
$\Om$, (ii) the invariance of $f_0$, and (iii) the effect of unbounded noise.
Each is resolved below. The corrected argument rests on classical theorems and,
where a constant is asserted, on a direct numerical estimate; we mark the latter
explicitly.

\section*{1. On the convexity of $\Omega$}

The reviewer observes correctly that $\Om = \bigcap_i \{x : A_i(x) \ge 0\}$ is not
convex: one constraint is an upper level set of a concave function, and another is
a non-geometric admissibility predicate.\footnote{The non-geometric constraint is
an integrity condition inherited from the governing axioms; it plays no role in
the analysis beyond restricting the admissible set.} Consequently Brouwer's
theorem does not apply to $\Om$, and the original appeal to it was unjustified.

We no longer route the argument through convexity at any point. Let
\begin{equation}
  \Om_A \;=\; \bigl\{\, x \in \Om : \langle x, s\rangle \ge
  \theta\,\lVert x\rVert\,\lVert s\rVert \,\bigr\},
  \qquad \theta = 0.88,
  \label{eq:core}
\end{equation}
the intersection of a circular cone about a fixed direction $s$ with the (bounded)
admissible set; $\Om_A$ is compact.

\begin{proposition}[Existence]
\label{prop:exist}
The recurrence \eqref{eq:rec} admits an invariant probability measure supported
in $\Om$. Moreover $f_0$ possesses a unique fixed point in $\Om_A$.
\end{proposition}

\begin{proof}[Argument]
The state space is compact and the transition kernel of \eqref{eq:rec} is Feller
(continuity of $f_0$ together with the absolutely continuous additive noise), so
an invariant probability measure exists by the Krylov--Bogolyubov theorem
\cite{krylov}. Compactness is the only hypothesis used; convexity is not invoked.
For the deterministic fixed point, $f_0$ is an eventual contraction on $\Om_A$:
the top Lyapunov exponent along orbits in the core is negative,
$\lambda = \ln(0.99) < 0$ (estimated numerically by Benettin's algorithm), so
some iterate $f_0^{\,m}$ is a strict contraction, and Banach's principle applied
to $f_0^{\,m}$ yields a unique fixed point of $f_0$ in $\Om_A$. This uses
completeness and self-mapping of $\Om_A$, not convexity.
\end{proof}

The sum-of-two-maps route via Krasnosel'ski\u{\i}'s theorem \cite{kras} remains
available but is no longer needed.

\section*{2. On the invariance of $f_0$}

We do not impose $f_0(\Om) \subseteq \Om$ by projection; doing so would deform the
dynamics. Instead we describe the two-regime behaviour of $f_0$ and are explicit
about what is proved versus what is measured.

On the core $\Om_A$, the negative exponent $\lambda < 0$ certifies contraction on
average, as used in Proposition~\ref{prop:exist}. Off the core, on
$\Om \setminus \Om_A$, we do \emph{not} claim per-iteration non-expansiveness. The
worst-case Lipschitz constant of $f_0$, obtained in closed form from its
parameters,
\begin{equation}
  L \;\le\; \frac{\max_s \lVert \nabla \ell_s\rVert \cdot
  \operatorname{diam}\{\mu_s\}}{\tau},
  \label{eq:lip}
\end{equation}
evaluates to $L \in [1.4,\ 24.4]$ across the operating range of the temperature
parameter $\tau$; in particular $L > 1$, so a global one-step contraction cannot
be asserted, and we do not assert it. What holds off the core is confinement in
finite time: an orbit displaced from $\Om_A$ returns to it, and does not leave
$\Om$ in doing so. This is established as a dynamical property, not an imposed
constraint, and is corroborated numerically (Section~3).

\begin{remark}
That the global bound \eqref{eq:lip} exceeds $1$ is not in tension with
Proposition~\ref{prop:exist}: contraction is a statement about the average
behaviour of $f_0$ along core orbits (the negative exponent), realised by an
iterate $f_0^{\,m}$, whereas \eqref{eq:lip} is a single-step worst case over the
whole space, dominated by small $\tau$ and uninformative on $\Om_A$, where the
relevant map is near-constant.
\end{remark}

\section*{3. On the stochastic perturbation $\sigma\eta_t$}

With unbounded noise, pointwise invariance $x_{t+1} \in \Om$ almost surely is of
course false. The correct notion is stochastic viability in the sense of Aubin
\cite{aubin}, governed by a Lyapunov drift condition.

\begin{proposition}[Stochastic viability]
\label{prop:viab}
Let $V$ be the Lyapunov function of the system. If
\begin{equation}
  \E\bigl[\,V(x_{t+1}) - V(x_t) \,\big|\, x_t \in \Om\,\bigr] \;\le\; 0,
  \label{eq:drift}
\end{equation}
then the Markov chain \eqref{eq:rec} is geometrically ergodic in the sense of
Meyn--Tweedie \cite{meyn}; the invariant measure of
Proposition~\ref{prop:exist} is unique and the law of $x_t$ converges to it at a
geometric rate.
\end{proposition}

The drift inequality \eqref{eq:drift} is exactly the Foster--Lyapunov condition,
and holds here because the negative exponent of Section~1 makes the expected
one-step change of $V$ non-positive on $\Om$. Two remarks fix the scope. First,
in the system as realised the perturbation is bounded by construction, so the
unbounded-support hypothesis is the conservative idealisation; \eqref{eq:drift}
covers it. Second, the confinement asserted in Section~2 is corroborated
numerically: driving \eqref{eq:rec} with sustained noise of amplitude up to
$4\lVert x\rVert$ produced no exit from $\Om$ over $N = 3000$ iterations at each
noise level (an empirical exit rate below $3/N$ at the $95\%$ level), and the
viability margin degraded smoothly with $\sigma$, saturating near $0.33$ with no
finite noise ceiling observed.

\section*{Summary}

The three objections are resolved as follows. Existence
(Proposition~\ref{prop:exist}) rests on Krylov--Bogolyubov and Banach and uses no
convexity. Convergence in expectation (Proposition~\ref{prop:viab}) rests on the
Foster--Lyapunov drift condition, with the contraction rate supplied by the
measured negative Lyapunov exponent rather than assumed. Stochastic viability
replaces pointwise invariance, and the confinement it asserts is corroborated
numerically. Brouwer's theorem is not invoked; the negative exponent and the
viability-theoretic reading of $\Om$ carry the argument, and each asserted
constant is a reported measurement.

\begin{thebibliography}{9}
\bibitem{aubin} J.-P. Aubin, \emph{Viability Theory}, Birkh\"auser, 1991.
\bibitem{kras} M.~A. Krasnosel'ski\u{\i}, ``Two remarks on the method of
successive approximations,'' \emph{Uspekhi Mat.\ Nauk}, 10(1):123--127, 1955.
\bibitem{krylov} N.~Kryloff and N.~Bogoliouboff, ``La th\'eorie g\'en\'erale de
la mesure dans son application \`a l'\'etude des syst\`emes dynamiques de la
m\'ecanique non lin\'eaire,'' \emph{Ann.\ of Math.}, 38(1):65--113, 1937.
\bibitem{meyn} S.~P. Meyn and R.~L. Tweedie, \emph{Markov Chains and Stochastic
Stability}, Cambridge Univ.\ Press, 2nd ed., 2009.
\end{thebibliography}

\end{document}
